Este guia orienta a configuração e utilização da ferramenta de otimização de forma direta e autônoma.

\subsection{Pré-requisitos do sistema}

Para executar esta ferramenta, seu ambiente deve atender a dois requisitos fundamentais:

\subsubsection{Python 3.10}
Verifique se o Python está instalado abrindo o \textit{Prompt de Comando} (cmd) e digitando:

\begin{verbatim}
py --list
\end{verbatim}

\textbf{Como interpretar:}
\begin{itemize}
    \item Se aparecer \texttt{ - 3.10.X} (ou superior), o Python está pronto.
    \item Se der erro, baixe o instalador em \texttt{python.org} e marque a opção "Add to PATH" durante a instalação.
\end{itemize}

\subsubsection{IBM ILOG CPLEX Optimization Studio}
O modelo utiliza a biblioteca \texttt{docplex}, que requer o motor de otimização da IBM instalado no sistema.

\begin{itemize}
    \item É obrigatório ter uma instalação do CPLEX com licença acadêmica ou comercial.
    \item Não utilize a versão "Community Edition", ela possui limites restritos de variáveis e restrições que são insuficientes para a complexidade deste case, impedindo a resolução.
\end{itemize}


\subsection{Configuração do ambiente}

Este passo é necessário apenas na primeira vez que você for usar a ferramenta. Localize e execute o arquivo \texttt{criar\_ambiente.bat} na pasta raiz do projeto.
\begin{itemize}
    \item Uma janela do terminal se abrirá e instalará automaticamente as dependências.
    \item Aguarde até ver a mensagem "AMBIENTE CRIADO COM SUCESSO" e feche a janela.
    \item Se a janela fechar instantaneamente sem mensagens de sucesso, verifique sua instalação do Python.
\end{itemize}

\subsection{Execução da ferramenta}

Para utilizar o otimizador no dia a dia:

\begin{enumerate}
    \item Execute o arquivo \texttt{executar\_solver.bat} na pasta raiz.
    \item Uma janela de seleção de arquivos abrirá automaticamente. Selecione a planilha Excel com os dados de entrada.
    \item Acompanhe o progresso no terminal. O solver processará os dados e buscará a melhor solução.
    \item Ao finalizar, a ferramenta abrirá automaticamente a planilha de resultados gerada na pasta \texttt{resultados}.
\end{enumerate}

\vspace{0.3cm}

\textbf{Controle de Execução (Importante):}

O solver foi configurado para buscar a solução ótima sem um limite de tempo pré-fixado.
\begin{itemize}
    \item \textbf{Interrupção manual:} Você pode interromper o processo a qualquer momento pressionando o atalho \texttt{Ctrl + C} na janela do terminal.
    \item \textbf{Salvamento automático:} Ao interromper o processo, o programa salvará automaticamente a melhor solução encontrada até aquele instante, garantindo que nenhum progresso válido seja perdido.
\end{itemize}


\subsection{Estrutura de arquivos}

A organização das pastas do projeto segue o padrão abaixo. Mantenha essa estrutura para garantir o funcionamento dos scripts automáticos.

\begin{verbatim}
desafio-otimizacao-suzano/
|
+-- criar_ambiente.bat          (Configuracao inicial)
+-- executar_solver.bat         (Execucao diaria)
+-- requirements.txt            (Lista de dependencias)
|
+-- code/                       (Codigos-fonte Python)
|   +-- modelo.py
|   +-- gestor_dados.py
|   +-- exportar_resultados.py
|
+-- venv/                       (Gerado automaticamente)
+-- resultados/                 (Gerado automaticamente)
\end{verbatim}

\subsection{Solução de problemas comuns}

\begin{itemize}
    \item \textbf{Erro "Python não encontrado" ou comando "py" inválido:} O Python não está instalado corretamente. Reinstale marcando a opção "Add to PATH".
    \item \textbf{Erro "Ambiente não encontrado":} A pasta \texttt{venv} pode ter sido deletada. Execute novamente o arquivo \texttt{criar\_ambiente.bat}.
    \item \textbf{Erro "Arquivo não encontrado":} Certifique-se de executar os arquivos \texttt{.bat} sempre a partir da pasta raiz do projeto, e não de dentro das subpastas.
\end{itemize}

\vspace{0.5cm}
\textbf{Nota de portabilidade:} Esta ferramenta é portável. Para usá-la em outro computador, basta copiar a pasta inteira do projeto e rodar o \texttt{criar\_ambiente.bat} no novo local.

\newpage